\documentclass[sigconf,nonacm]{acmart}
\settopmatter{printfolios=true}

\usepackage{booktabs}
\usepackage{tabularx}
\usepackage[english,russian]{babel}

% Русские названия окружений теорем: acmart определяет свои только для
% неопределённых имён (\@ifundefined в \AtEndPreamble), поэтому объявляем их здесь.
\theoremstyle{acmplain}

\newcolumntype{Y}{>{\raggedright\arraybackslash}X}
\newtheorem{theorem}{Теорема}[section]
\newtheorem{proposition}[theorem]{Утверждение}
\newtheorem{lemma}[theorem]{Лемма}
\newtheorem{corollary}[theorem]{Следствие}
\theoremstyle{acmdefinition}
\newtheorem{example}[theorem]{Пример}
\newtheorem{definition}[theorem]{Определение}
\theoremstyle{acmplain}

% Long tokens (tool invocations, file names) make two-column justification
% tight; a little emergency stretch is preferable to overfull lines.
\emergencystretch=2em

% Inline code. Carets, percent signs and tildes are escaped at the call site.
\newcommand{\code}[1]{\texttt{#1}}
% Guarantee names.
\newcommand{\G}[1]{G#1}

\newtheorem{remarknum}{Замечание}

\begin{document}

\title{Выпуск полиглотных проектов как распределённая транзакция:\\
       протокол CCME и его реализация в dispat}

\author{Семён Федюкович}
\affiliation{%
  \institution{Независимый исследователь}
  \city{}
  \country{}}
\email{yohimik@gmail.com}

\renewcommand{\shortauthors}{Семён Федюкович}

\begin{abstract}
Выпуск взаимозависимых пакетов необратимо пишет в независимые реестры без
общего координатора и отката. Частичный отказ поэтому может навсегда оставить
потребителя позади опубликованной зависимости.

Статья представляет CCME, язык сообщений коммитов и протокол возобновляемого
выпуска, и dispat, его открытую реализацию для репозиториев с одним пакетом,
монорепозиториев и полирепозиториев, объединённых управляющим репозиторием. Состояние
выпуска хранится в метках git, план есть чистая функция истории, графа и
конфигурации, а распространённая работа допускается по окну самого потребителя.
Доказаны восемь гарантий, включая отсутствие осиротевших пакетов, однократное
погашение в журнале и сходимость при явном инварианте повторного запуска.
CCME~2.0 включает шесть исправлений, полученных из контрпримеров к исторической
спецификации~1.0, в том числе разные окна серии и свежей работы и допуск с
учётом каналов.

Свежая кампания на официальном образе dispat~1.8.2 прошла шесть выбранных
регрессий CCME~2.0, совпала с историческим оракулом~1.0 на всех 1{,}500
экземплярах core, завершила 24/24 эпизода с отказами не более чем за четыре
запуска и все восемь ячеек масштаба. Полный профиль намеренно создаёт
регрессивные предварительные версии: 997 планов приняты, а 503 полных
диагностических плана отклонены с E195. Все 1{,}914 различающихся строк
сохраняют эвристическую классификацию и указанные ограничения покрытия.
\end{abstract}

\begin{CCSXML}
<ccs2012>
<concept>
<concept_id>10011007.10011074.10011099.10011102</concept_id>
<concept_desc>Software and its engineering~Software configuration management and version control systems</concept_desc>
<concept_significance>500</concept_significance>
</concept>
<concept>
<concept_id>10002951.10002952.10003190</concept_id>
<concept_desc>Information systems~Database transaction processing</concept_desc>
<concept_significance>300</concept_significance>
</concept>
</ccs2012>
\end{CCSXML}
\ccsdesc[500]{Software and its engineering~Software configuration management and version control systems}
\ccsdesc[300]{Information systems~Database transaction processing}

\keywords{инженерия выпусков, монорепозиторий, полирепозиторий, семантическое
версионирование, Conventional Commits, саги, частичный отказ,
идемпотентность, детерминированное повторное исполнение}

\maketitle

\noindent\textbf{Снимок спецификации.} Эта статья анализирует выпущенную
спецификацию CCME~2.0.0~\cite{ccme2}. Если приводится неверное правило CCME~1.0.0 из
фиксации \code{713f1a0b}, оно и результат явно помечены как исторические.
CCME~3.0.0 --- историческое первое раскрытие будущей конструкции VCS и отката.
Выпущенная CCME~3.0.1~\cite{ccme301} исправляет рекомендации по алгоритму.
Dispat~1.8.2 не реализует ни один выпуск 3.0; модель и свидетельства E1--E5
описывают только реализованную семантику CCME~2.0.

\section{Введение}
\label{sec:intro}

Рабочее пространство выпуска может содержать один пакет, несколько пакетов
монорепозитория или пакеты из раздельно сопровождаемых репозиториев
(полирепозиторий). Выпуск означает выбор версии для пакетов, которым она нужна,
запись этих версий в манифесты,
сборку артефактов, их загрузку в реестры и протоколирование
произошедшего. В современном проекте реестров несколько, и они
разнородны: пакет npm соседствует с модулем Go, образом Docker и сайтом
документации, и все они связаны в один граф зависимостей.

Трудность составляет протоколирование. Независимые реестры не дают общей
атомарной транзакции: каждая публикация может завершиться отдельно, а её исход
может потеряться до записи соответствующей метки git. Это долгоживущая
транзакция над раздельно фиксируемыми ресурсами, постановка саг~\cite{sagas};
восстановление должно завершить или сверить успешные шаги, а не предполагать
межреестровый откат.

Lerna, Nx~\cite{nx}, release-please и Changesets~\cite{changesets}
автоматизируют выпуск, а semantic-release~\cite{semanticrelease} рассчитан на
один пакет в репозитории. Зафиксированные ниже эксперименты показывают разные
границы восстановления у реестровых запросов, файлов намерений и времени
расстановки меток.

Настоящая статья рассматривает выпуск как транзакцию, которой он и
является, и вносит три вклада:
\begin{itemize}
  \item язык сообщений фиксаций CCME, в котором намерение выпуска
    (область воздействия по двум отдельным, односторонне связанным осям,
    инкременты версий и
    предрелизные каналы) есть доступный для рецензирования текст в
    самой фиксации (раздел~\ref{sec:ccme});
  \item формальную модель вычисления выпуска и протокола публикации,
    самодостаточно изложенную в настоящей статье, с восемью доказанными
    над ней гарантиями (разделы~\ref{sec:model} и~\ref{sec:guarantees});
    формализация приводит исторические контрпримеры к спецификации~1.0
    и исправленные правила, вошедшие в CCME~2.0
    (раздел~\ref{sec:eval:formal});
  \item dispat, независимую от экосистемы реализацию, в которой пакет ---
    каталог, а стадия --- команда оболочки, оценённую дифференциальной
    проверкой, выбранными векторами CCME~2.0, кампаниями восстановления,
    воспроизведением и зафиксированными экспериментами с отказами
    (разделы~\ref{sec:impl} и~\ref{sec:eval}).
\end{itemize}

\section{Мотивирующий пример}
\label{sec:motivation}

Рабочее пространство содержит библиотеку \code{core} и три потребителя
\code{cli}, \code{ui}, \code{api}, каждый из которых зависит от
\code{core}. Поступает фиксация $C$ с сообщением
\code{feat(core)\^{}:~streaming reader}: \code{core} получает младший
инкремент (minor), а знак каре (указатель глубины распространения)
распространяет исправительный инкремент (patch) на одно ребро наружу.
Запуск публикует \code{core@1.5.0}, \code{ui@0.9.2} и \code{api@1.2.1},
помечая каждый при успехе, после чего публикация \code{cli} отказывает
по любой из причин, по которым отказывают записи в реестры.

Естественно ожидать, что повторный запуск доделает работу. Так ли это,
зависит от одного определения. Предположим, что механизм выпуска
сохраняет намерение выпуска фиксации в силе лишь до тех пор, пока оно
ещё ожидает выпуска в \emph{собственном} пакете фиксации; это
естественное прочтение. Тогда на втором запуске \code{core} уже выпустил
$C$, поэтому у семантической единицы нет пакетов-источников, она ни на
что не распространяется, и \code{cli} не выпускается никогда, ни на
этом, ни на любом последующем запуске. Пакет \code{cli}
\emph{осиротел}: он навсегда отстаёт на один исправительный инкремент от
объявленной им зависимости, без какой-либо диагностики, поскольку с
точки зрения механизма ничего не ожидает выпуска. Отказ незаметен и
становится тем вероятнее, чем больше разветвление графа.

Исправление есть изменение определения в одну строку: единица
распространяется на потребителя ровно до тех пор, пока окно ожидания
\emph{потребителя} ещё содержит фиксацию единицы
(определение~\ref{def:adm}). Догоняющий выпуск тогда не требует отдельного прохода или побочного состояния,
когда предыдущие исходы подтверждены метками; неопределённый разрыв между
внешней публикацией и меткой всё ещё требует сверки идентичности.
Раздел~\ref{sec:guarantees} доказывает гарантию, которую это даёт
(\G{2}, отсутствие осиротевших пакетов), и родственные ей гарантии.

\section{CCME: язык}
\label{sec:ccme}

CCME --- строгое надмножество Conventional Commits~1.0.0~\cite{cc} над
заявленным соответствующим подмножеством. Одно сообщение фиксации
представляет собой последовательность \emph{семантических единиц},
разделённых разделителем, занимающим целую строку (\code{---}); каждая
единица --- независимый блок в форме Conventional Commits со своими
заголовком, телом и служебными строками, так что запрос на слияние,
влитый со сжатием фиксаций, описывает несколько несвязанных изменений,
не приписывая всем им один ложный тип:

\begin{small}
\begin{verbatim}
feat(@acme/core)^^minor%beta++*: reader
---
fix(@acme/cli): correct exit code on SIGINT
---
cancel(@acme/legacy-adapter): reset state
\end{verbatim}
\end{small}

\emph{Области действия именуют пакеты или выводятся.} Область действия
задаётся множеством пакетов через запятую или шаблонами с исключениями
(\code{feat(*,-docs)}) либо отсутствует; в последнем случае затронутое
множество выводится из файлов, изменённых фиксацией, с определением
пакета по каталогу-владельцу. Именно вывод заставляет обычные фиксации
без области действия выпускать ровно то, чего они коснулись.

\emph{Две оси, две глубины.} Фиксация отвечает на два отдельных вопроса
о своих потребителях: как далеко распространяется \emph{инкремент
версии} и как далеко распространяется \emph{предрелизный канал}. По
умолчанию обе глубины равны нулю: фиксация, которая ничего не говорит,
выпускает собственные пакеты и ничего больше. Ось инкрементов
записывается как \code{\^{}} (прямые потребители), \code{\^{}\^{}}
(транзитивное замыкание) или \code{+N}; ось каналов --- как \code{\%}
(собственные пакеты меняют канал), \code{\%\%} (меняют и потребители),
\code{++N} (насколько далеко). Каждый знак-модификатор раскрывается в
служебную строку; приведённый выше заголовок раскрывается в
\begin{small}
\begin{verbatim}
feat(@acme/core): reader

Propagate: minor          Channel: beta
Propagate-Depth: all      Propagate-Channel-Depth: all
\end{verbatim}
\end{small}
что читается так: \code{core} получает младший инкремент и переходит в
серию beta, каждый транзитивный потребитель получает младший инкремент и
следует за ним на \code{beta}. Значения канала могут быть
\emph{переходами} \code{<from>><to>}: переместить только пакеты,
находящиеся сейчас на \code{<from>}. Переходы делают перевод в
стабильный канал идемпотентным:
\code{release(@acme/*)\%beta>stable:~ship} переводит в стабильный канал
всё, что ещё остаётся на beta, и может быть безопасно исполнен повторно
после частичного отказа.

\begin{table}[t]
\caption{Примеры заголовков.}
\label{tab:headers}
\small
\begin{tabular}{@{}ll@{}}
\toprule
Заголовок & Значение \\
\midrule
\code{feat: add retry} & выводимая область, minor \\
\code{fix(api): null guard} & \code{api}, patch \\
\code{perf(core)\^{}patch+1: x} & patch прямым потребителям \\
\code{refactor(core)\^{}\^{}major!: x} & несовместимое; major замыканию \\
\code{feat(cli)\%beta: watch mode} & \code{cli} в серию beta \\
\code{feat(core)\^{}\%beta++1: x} & сам на beta; потребители следом \\
\code{release(a)\%beta>stable: ship} & в stable оставшиеся на beta \\
\code{cancel(*): reset state} & отбросить ожидающие записи \\
\bottomrule
\end{tabular}
\end{table}

\emph{Отмена и корректирующие указания.} Тип \code{cancel} отбрасывает
накопленные невыпущенные метаданные выпуска для области действия;
единица отменяется для пакета, когда её фиксация есть предшественник
отменяющей фиксации или совпадает с ней; отношение предшествования, в
отличие от отметок времени, детерминировано при слияниях и
перебазированиях. Служебные строки \code{Edits:~<sha>} и
\code{Deletes:~<sha>} переформулируют или отбрасывают названную прошлую
запись: фиксация \code{fix(core):~rewrite internals}, несущая
\code{Edits:~<sha of A>}, заменяет ожидающий старший инкремент фиксации
$A$ исправительным инкрементом этой единицы, так что \code{core}
выпускает \code{1.4.3}, а не \code{2.0.0}. Корректирующие указания
достигают только невыпущенной работы и только строгих предшественников;
побеждает новейшее указание; указание, чья собственная запись отброшена,
ничтожно. \code{Release-As:~none} удерживает пакет до снятия удержания;
точное значение \code{Release-As} закрепляет одну версию.

\emph{Диагностика и разбор.} Перечень из 94 кодов диагностик (40 ошибок,
54 предупреждения) закрепляет и исходы, и объяснения, так что две
соответствующие реализации не расходятся ни в том, ни в другом; восемь
предупреждений неподавляемы, поскольку каждое объясняет запись плана,
которую читатель одного лишь журнала объяснить не может (догоняющий
выпуск, заблокированный пакет, подавленное распространение). Поскольку
сообщение есть недоверенный ввод, разбираемый в системе непрерывной
интеграции, соответствующий синтаксический анализатор представляет
собой ограниченное однопроходное сканирование за время $O(n)$ с рабочей
памятью $O(1)$ и без возвратов, так что враждебное сообщение не может
вызвать сверхлинейного разбора.
\section{Модель}
\label{sec:model}

Этот раздел определяет вычисление выпуска и протокол публикации;
раздел~\ref{sec:guarantees} доказывает гарантии над этими
определениями и ни над чем иным.

\subsection{Рабочее пространство, состояние, окна}

\begin{definition}[история]
История есть конечный ациклический граф фиксаций с выделенной вершиной ветви
$\mathsf{HEAD} = h$; $\mathit{reach}(c)$ есть множество фиксаций, достижимых
из $c$ по рёбрам родительства, включая саму $c$; запись $c' \preceq c$ означает
$c' \in \mathit{reach}(c)$. Вершина $h$ неизменна на всём протяжении изложения. Такое
прочтение репозитория стандартно: ациклический граф фиксаций с отношением достижимости служит
рабочей моделью истории git в анализе репозиториев~\cite{birdgit}, и
та же структура лежит в основе формальных выводов трёхстороннего
слияния git из наименьших общих предшественников~\cite{mrdt}. В отличие от
описаний управления версиями на основе теории патчей
(patch theory)~\cite{swierstraloh,hpt,mimram}, данная модель вовсе не отслеживает
содержимое файлов, а только историю, метки и сообщения, что и
делает план чистой функцией репозитория.
\end{definition}

\begin{definition}[рабочее пространство]
$\mathcal{P}$ есть конечное множество пакетов в $h$; $E \subseteq
\mathcal{P} \times \mathcal{P}$ содержит $(a,b)$, когда $a$ объявляет
зависимость от $b$ сконфигурированного вида; $\mathit{dep}(p) = \{d : (d,p)
\in E\}$. Граф ацикличен (цикл есть ошибка уровня репозитория: для него
не существует порядка публикации, так что любой участник цикла, опубликованный первым,
объявил бы диапазон на несуществующую версию). Для $S \subseteq
\mathcal{P}$ величина $\delta(S,d)$ есть наименьшее число рёбер на цепи от
элемента $S$ к $d$ против направления стрелок и $\infty$, если такой цепи нет.
\end{definition}

\begin{definition}[инкременты, версии, каналы]
$\mathbb{B} = \{\mathsf{none} < \mathsf{patch} < \mathsf{minor} <
\mathsf{major}\}$ есть цепь, а потому верхняя полурешётка относительно
$\sqcup = \max$. Версии следуют SemVer~2.0.0~\cite{semver} с его отношением
старшинства; $\mathit{ch}(v)$ равно $\mathsf{stable}$ для непредрелизной
$v$, а иначе равно первому предрелизному идентификатору $v$; $\mathit{core}(v)$
и $\mathit{cnt}(v)$ суть числовая тройка и предрелизный счётчик;
$\mathit{apply}(v,b)$ применяет инкремент.
\end{definition}

\begin{definition}[состояние]
\label{def:state}
Состояние $\sigma$ есть конечное множество меток выпуска $(p, v, c)$, где
$c \preceq h$; состояния упорядочены по включению. $\mathit{base}_\sigma(p)$
есть версия наибольшего старшинства среди меток $p$,
$\mathit{sbase}_\sigma(p)$ есть \emph{непредрелизная} версия наибольшего старшинства,
$\mathit{basec}_\sigma(p)$ и $\mathit{scom}_\sigma(p)$ суть их
фиксации; по соглашению $\mathit{ch}(\mathit{base}_\sigma(p)) = \mathsf{stable}$,
когда у $p$ нет ни одной метки. Метки никогда не перемещаются и не удаляются, поэтому
состояния только растут: метка есть запись о том, что шаг транзакции
завершён, а переписанная запись и есть то, что делает сагу невосстановимой.
Поэтому право записи меток и есть полномочие на выпуск, а манифесты никогда не
читаются как состояние, так что подложный манифест в ответвлённом репозитории (fork) не может
повлиять на версионирование.
\end{definition}

\begin{definition}[окна серии и свежей работы]
\label{def:window}
\begin{equation}
  W_\sigma(P) \;=\; \mathit{reach}(h) \setminus
                    \mathit{reach}(\mathit{scom}_\sigma(P)),
  \label{eq:window}
\end{equation}
и
\begin{equation}
  W^\circ_\sigma(P)=W_\sigma(P)\setminus
  \mathit{reach}(\mathit{basec}_\sigma(P)), \label{eq:fresh}
\end{equation}
где при отсутствии любой метки вычитаемое множество пусто. $W_\sigma$ есть
окно серии, а $W^\circ_\sigma$ --- окно свежей, ещё не доставленной работы.
$W_\sigma(P) = \mathit{reach}(h)$, когда у $P$ нет стабильной базовой метки:
первый выпуск ещё не выпускавшегося пакета учитывает всю его
историю, и именно поэтому внедрение на существующем репозитории не требует
шага миграции. Окно отсчитывается от последней \emph{стабильной} метки,
чтобы предрелизная серия могла заново вычислять свою целевую версию по всей
серии (уравнение~\ref{eq:train}). Прямой допуск, допуск распространения и
погашение реестра используют $W^\circ$: предрелизная метка погашает уже
доставленную работу, не стирая данные, нужные для цели всей серии.
\end{definition}

\subsection{Единицы, подавление, допуск}

\begin{definition}[единицы]
$\mathcal{U}$ есть множество семантических единиц, переносимых фиксациями из
$\bigcup_P W_\sigma(P)$, после разбора, с разрешёнными областями действия и применёнными
корректирующими указаниями. Для $u \in \mathcal{U}$: $c_u$ есть её фиксация,
$\mathit{scope}(u) \subseteq \mathcal{P}$ её разрешённая область действия,
$\mathit{bump}(u) \in \mathbb{B}$ инкремент, в который отображаются её тип и \code{!};
её ось инкрементов составляют $\mathit{pb}(u) \in \mathbb{B}$ (указанное значение либо
$\mathit{bump}(u)$ при \code{inherit}; значение по умолчанию, как только запрошено
распространение, есть $\mathsf{patch}$, честный сигнал для потребителя, чей
собственный API не изменился), $\mathit{pd}(u) \in \mathbb{N} \cup
\{\infty\}$ (по умолчанию $0$) и фильтр области действия $\mathit{ps}(u)$; её ось
каналов составляют значение канала $\mathit{pc}(u)$, глубина $\mathit{pcd}(u)$
(по умолчанию $0$) и фильтр $\mathit{pcs}(u)$.
\end{definition}

\begin{definition}[подавление и источники]
\label{def:src}
$\mathit{canc}(c,p)$ выполняется, когда $c$ есть предшественник фиксации
\code{cancel} (или сама эта фиксация), область действия которой содержит $p$; $\mathit{held}(p)$
выполняется, когда действующее для $p$ значение \code{Release-As} равно \code{none}. Оба предиката
суть функции только $h$ и конфигурации. Полагая $\mathit{disch}_\sigma(p,c)
\equiv c \notin W^\circ_\sigma(p)$, имеем
\begin{equation}
  \mathit{src}_\sigma(u) = \{ p \in \mathit{scope}(u) :
    \mathit{disch}_\sigma(p,c_u) \lor
    \lnot(\mathit{canc}(c_u,p) \lor \mathit{held}(p)) \}.
  \label{eq:src}
\end{equation}
Подавление применяется только к \emph{непогашенной} работе: как только $p$
опубликовал версию, несущую $u$, артефакт, причитающийся его потребителям,
стал публичным, и ни последующая отмена, ни последующее удержание не вправе задним числом
лишить их этого артефакта.
\end{definition}

\begin{definition}[достижимое множество]
\label{def:reach}
$R(S,\ell) = \{ d \in \mathcal{P} \setminus S : \delta(S,d) \leq \ell
\}$. Исключение всего $S$ означает, что единица никогда не распространяется на пакет,
который она инкрементирует напрямую; глубина измеряется по кратчайшему пути и всегда
от исходного множества источников, а распространённый эффект никогда не порождает
каскада собственного распространения, поэтому глубина, записанная в сообщении, есть вся
область воздействия в каждом запуске.
\end{definition}

\begin{definition}[каналы]
\label{def:chan}
$\mathit{chan}_\sigma(P)$ разрешается для \emph{каждого} пакета: это значение,
предложенное новейшей прямой директивой \code{Channel}, называющей $P$, с
$c_u \in W^\circ_\sigma(P)$; иначе значение, предложенное для $P$ новейшей
единицей $u$ оси каналов с $\mathit{pcd}(u) > 0$, $P \in
R(\mathit{src}_\sigma(u), \mathit{pcd}(u)) \cap \mathit{pcs}(u)$,
$c_u \in W^\circ_\sigma(P)$ и $\lnot\mathit{canc}(c_u,P)$; иначе
$\mathit{ch}(\mathit{base}_\sigma(P))$. Предложение вычисляется относительно
$\mathit{cur} = \mathit{ch}(\mathit{base}_\sigma(P))$ (метки, но никогда не
значения, вычисленного в этом запуске) и пусто, когда оно предлагает
сам $\mathit{cur}$; переход предлагает значение лишь там, где $\mathit{cur}$
совпадает с его \code{<from>}; \emph{распространённое} значение $\mathsf{stable}$,
не являющееся переходом, никогда не переводит предрелизный пакет в стабильный канал (такой
перевод должен быть намеренным); \code{inherit} разрешается в собственное
значение \code{Channel} единицы, если оно есть, а иначе в
$\mathit{ch}(\mathit{base}_\sigma(q))$ для её источников $q$, причём разногласие
решает побайтово наименьшее имя.
\end{definition}

\begin{definition}[допуск и результирующий инкремент]
\label{def:adm}
Пусть $\mathit{live}(u)$ означает $\mathit{bump}(u)\neq\mathsf{none}$,
$\mathit{pb}(u)\neq\mathsf{none}$ и $\mathit{pd}(u)>0$. Пусть $\mathit{srcCh}_\sigma(u)$ равно
$\{\mathit{chan}_\sigma(p) : p \in \mathit{src}_\sigma(u)\}$. Тогда при
\begin{equation}
  \mathit{rsv}_\sigma(u,D) \equiv \mathsf{stable} \in
  \mathit{srcCh}_\sigma(u) \lor \mathit{chan}_\sigma(D) \in
  \mathit{srcCh}_\sigma(u)
  \label{eq:resolvable}
\end{equation}
(установщик выбирает предрелизную версию только на её собственной линии, поэтому инкремент
через границу каналов не является действительным обязательством),
\begin{equation}
\label{eq:admission}
\begin{aligned}
  \mathit{adm}_\sigma(u,D) \;\equiv\;\;
     & D \in R(\mathit{src}_\sigma(u), \mathit{pd}(u))
     &\!\!& \text{(достигнут)} \\
  \wedge\;\; & D \in \mathit{ps}(u)
     &\!\!& \text{(в области действия)} \\
  \wedge\;\; & c_u \in W^\circ_\sigma(D)
     &\!\!& \text{(окно \textbf{потребителя})} \\
  \wedge\;\; & \neg\,\mathit{canc}(c_u, D)
     &\!\!& \text{(не отброшен)} \\
  \wedge\;\; & \mathit{live}(u)
     &\!\!& \text{(живой инкремент)} \\
  \wedge\;\; & \mathit{rsv}_\sigma(u, D)
     &\!\!& \text{(разрешим),}
\end{aligned}
\end{equation}
\begin{equation}
\label{eq:effective}
  \mathit{eff}_\sigma(P) = \mathit{dir}_\sigma(P) \sqcup
  \bigsqcup \{\, \mathit{pb}(u) : \mathit{adm}_\sigma(u,P) \,\},
\end{equation}
где $\mathit{dir}_\sigma(P) = \bigsqcup\{\mathit{bump}(u) : P \in
\mathit{src}_\sigma(u) \cap \mathit{scope}(u),\ c_u \in W^\circ_\sigma(P)\}$.
Поскольку $\sqcup$ есть точная верхняя грань, на пакет с собственным ожидающим \code{feat}
входящий распространённый \code{patch} не влияет, распространение может
только повысить версию, а порядок сбора не имеет значения.
\end{definition}

Проверка устаревания со стороны потребителя есть восходящая двойственность
формулы~\ref{eq:admission}, а не только обратная достижимость. Для $D$ она
просматривает единицы в $W^\circ_\sigma(D)$, исключает $D$, если он сам входит
в источники единицы, и применяет те же кратчайшую глубину, область, отмену, $\mathit{live}$ и
предикат каналов $\mathit{rsv}$. Поэтому восходящий ответ точно совпадает с
обязательствами нисходящего планировщика, в том числе для предрелизов.

\subsection{План, версии, запуски}

\begin{definition}[план и версии]
\label{def:plan}
Пусть $\mathit{pin}_\sigma(P)=v$, если действующая сохранившаяся точная единица
\code{Release-As: v} для $P$ имеет $c_u\in W^\circ_\sigma(P)$; иначе
$\mathit{pin}_\sigma(P)$ не определено. Успешная метка погашает закрепление,
как и прочую свежую работу, а директива неизменна при повторной попытке.
\begin{equation}
\begin{aligned}
  \Pi(\sigma) = \{\, P :\;
    &\lnot\mathit{held}(P) \wedge (\mathit{eff}_\sigma(P) \neq
     \mathsf{none} \\[-1mm]
    &\ \lor\ \mathit{chan}_\sigma(P) \neq
     \mathit{ch}(\mathit{base}_\sigma(P))
     \ \lor\ \mathit{pin}_\sigma(P)\downarrow) \,\},
\end{aligned}
\end{equation}
причём каждый элемент несёт $\nu_\sigma(P)$. Пусть
$S=\mathit{sbase}_\sigma(P)$, $B=\mathit{base}_\sigma(P)$, а $E$ ---
совокупный инкремент серии по $W_\sigma(P)$. Если
$\mathit{pin}_\sigma(P)\downarrow$, оно задаёт точную версию. Стабильный пакет, остающийся стабильным,
использует $\mathit{apply}(S,\mathit{eff}_\sigma(P))$; предрелизный пакет,
переходящий в stable, использует $\mathit{apply}(S,E)$ по \S11.5 CCME~2.0.
Иначе, при целевом канале $\mathit{ch}$,
\begin{equation}
\label{eq:train}
\begin{aligned}
  t &= \mathit{apply}(S, E), \\
  \nu_\sigma(P) &=
  \begin{cases}
    t\text{-}\mathit{ch}.(\mathit{cnt}(B){+}1)
      & \text{если } \mathit{core}(B) = t \wedge \mathit{ch}(B) =
        \mathit{ch},\\
    t\text{-}\mathit{ch}.0 & \text{иначе,}
  \end{cases}
\end{aligned}
\end{equation}
с одним защитным условием: каждая вычисленная или закреплённая версия обязана превосходить
$\mathit{base}_\sigma(P)$ по старшинству, в противном случае запуск отказывает. Именно повторное
вычисление $t$ от \emph{стабильной} базовой метки в каждом запуске переводит всю
серию на \code{2.0.0-beta.0}, когда несовместимое изменение попадает в середину серии
(таблица~\ref{tab:train}).

Явное правило выпуска из серии в \S11.5 задаёт этот совокупный инкремент;
резюме \S13.9 вместо этого применяет свежий $\mathit{eff}$ к $S$ и должно
ссылаться на \S11.5. Формулировка E185 также требует, чтобы цель строго
превосходила ядро базовой версии, хотя нормативный пример переводит
\code{1.3.0-rc.1} в равное ядро \code{1.3.0}; защита должна отличать
регрессию от обычного перехода к равному ядру. Это несогласованности ссылок
и формулировки защиты, а не утверждение о потере цели реализацией.
\end{definition}

\begin{table}[t]
\caption{Уравнение~\ref{eq:train}, прослеженное на одной серии.}
\label{tab:train}
\small
\begin{tabularx}{\columnwidth}{@{}lYY@{}}
\toprule
Базовая метка & Ожидающий инкремент, канал & Следующая версия \\
\midrule
\code{1.2.3}        & младший (minor), \code{beta}     & \code{1.3.0-beta.0} \\
\code{1.3.0-beta.0} & нет свежей фиксации              & пустой план \\
\code{1.3.0-beta.0} & свежая фиксация minor            & \code{1.3.0-beta.1} \\
\code{1.3.0-beta.1} & младший $+$ новое исправление    & \code{1.3.0-beta.2} \\
\code{1.3.0-beta.2} & приходит несовместимое изменение & \code{2.0.0-beta.0} \\
\code{1.3.0-beta.2} & нет, переход на \code{rc}        & \code{1.3.0-rc.0} \\
\code{1.3.0-rc.1}   & нет, \code{stable}               & \code{1.3.0} \\
\bottomrule
\end{tabularx}
\end{table}

\begin{definition}[запуск]
\label{def:run}
Запуск из $\sigma$ параметризован выбираемым противником множеством
$F \subseteq \Pi(\sigma)$ отказавших публикаций, которое индуцирует
заблокированное множество
$B = \{p \in \Pi(\sigma) \setminus F : p$ транзитивно зависит от некоторого
$q \in F\}$, вычисляемое по \emph{полному} графу рабочего пространства (пакет
без инкремента в этом запуске всё равно образует путь от потребителя к отказавшему
поставщику). Публикация выполняется в топологическом порядке полного
графа, равенства разрешаются побайтовым сравнением имён, с фильтрацией по плану \emph{в последнюю очередь}:
если сначала индуцировать граф на план, исчезнут упорядочивающие пути, проходящие
через незапланированные пакеты. Каждый пакет получает метку сразу после успеха его
собственной публикации, никогда раньше (метка для неопубликованной версии
делает её невосстановимой, поскольку уравнение~\ref{eq:window} считает её
выпущенной), никогда группой по завершении (отказ посреди запуска мог бы
повторить внешнюю запись, которую цель может отвергнуть или выполнить дважды) и \emph{независимо от цели публикации},
включая цели, которые ничего не загружают
(утверждение~\ref{prop:tagall}). Запуск даёт
$\mathit{run}(\sigma,F) = \sigma \cup \{(p, \nu_\sigma(p), h) : p \in
\Pi(\sigma) \setminus (F \cup B)\}$. \emph{Последовательность запусков} есть
$\sigma_1 \leq \sigma_2 \leq \cdots$ с $\sigma_{k+1} =
\mathit{run}(\sigma_k, F_k)$.
\end{definition}

Для оценки сложности пусть $H,A$ считают достижимые фиксации и родительские
рёбра, $C$ --- фиксации объединённого окна, $U$ --- единицы, $I$ ---
разрешённые связи единица--пакет, а $Z$ --- сохранённые или выведенные связи
вкладов и происхождения. Классы окон могут совместно использовать неизменяемые
множества достижимости, но их построение стоит в $H+A$, а не только в $C$;
стабильные и свежие окна имеют разные ключи границ. Буквальный обход
распространения стоит $O(U(|\mathcal P|+|E|)+Z)$. Обход графа можно кэшировать
безопасно кэшировать при совпадении множества источников, глубины и видов рёбер;
иные оптимизации требуют отдельного доказательства эквивалентности; объединение по
окну и написанию директив небезопасно, поскольку допуск, отмена, область,
исключение собственного источника, наследуемый канал и происхождение остаются
свойствами единицы. Лексический разбор и разрешение областей имеют разные ключи
кэша. Обращение связей канала стоит $O(U+I+|\mathcal P|)$. Это информативные
оценки реализации, а не обязательные сквозные границы; нормативны
семантическая эквивалентность и детерминированный вывод, включая происхождение.
Векторы CCME~3.0.1 82h--82j проверяют разрешение прямого канала,
происхождение наследуемого канала и допуск отдельной единицы, не меняя G1--G8 и
реализованный алгоритм CCME~2.0.

Порядок важен по причине, которую метки выразить не могут: потребитель, опубликованный
раньше своего поставщика, несёт диапазон, согласованный со старой
версией, оба попадают в одну и ту же фиксацию, и никакой последующий запуск не может обнаружить
эту инверсию по отношению предшествования, поэтому она предотвращается, а не диагностируется. Когда
публикация отклонена из-за того, что версия уже существует (запуск погиб
между публикацией и меткой), CCME~2.0 требует сверить тождество содержимого:
тождественный артефакт принимается с предупреждением, несовпадение даёт
ошибку. Это нормативное требование не означает, что ядро dispat умеет
проверять произвольный сценарий публикации: интеграция должна предоставить
идентификатор содержимого или идемпотентную операцию.

\begin{theorem}[детерминизм]
\label{thm:det}
$\Pi(\sigma)$, $\nu_\sigma$, порядок публикации и диагностики суть
функции только от $(\sigma, h, \mathit{config})$.
\end{theorem}
\begin{proof}
Каждое определение выше есть замкнутое выражение от этих входов; три
точки выбора (конкурирующие директивы, конкурирующие распространённые каналы,
равенства в порядке публикации) разрешаются каждая явно заданным линейным порядком: новейшая
фиксация, затем последняя единица внутри неё, и побайтово наименьшее имя.
\end{proof}

\subsection{Конкретизация: что теоремы говорят о репозитории}
\label{sec:model:inst}

Приведённые выше определения параметричны, и теорема о модели
ничего не стоит для оператора, если это не теорема о
\emph{его} репозитории. Конкретный репозиторий в конкретной фиксации определяет
один \emph{экземпляр}
\[
  I \;=\; (\text{фиксации и рёбра родительства},\ h,\
           \mathcal{P}, E,\ \mathcal{U},\ \sigma,\ \mathit{config}):
\]
ациклический граф фиксаций и \code{HEAD} читаются из репозитория,
пакеты и рёбра из его манифестов и конфигурации, единицы
из его сообщений фиксаций, а $\sigma$ из его меток. Каждая величина,
определяемая моделью, есть вычислимая функция от $I$: множества
конечны, отношение предшествования разрешимо, и ни одно определение не обращается ни к чему
вне $I$. Модель описывает фрагмент поведения репозитория: она абстрагирует
содержимое артефактов, семантику перехватчиков и не представленные в $I$
возможности конфигурации. Теоремы раздела~\ref{sec:guarantees} универсальны
для экземпляров этого фрагмента при ацикличности, защитном условии на версию
из определения~\ref{def:plan} и сформулированном ниже инварианте повторного
запуска. Последний есть свойство последовательности запусков, проверяемое
перед каждой попыткой, а не статическое свойство одного экземпляра. Между теоремами и работающим исполняемым файлом остаётся
утверждение об \emph{уточнении}: что реализация вычисляет
$\Pi(\sigma)$ и $\nu_\sigma$ этой модели и исполняет
протокол определения~\ref{def:run}. Это утверждение здесь не доказывается;
оно проверяется тестами: проверочными векторами раздела~\ref{sec:eval:fresh} (каждый
из них задаёт небольшой экземпляр и план, который модель ему назначает) и
разделом~\ref{sec:eval:diff} на случайно порождённых экземплярах.
\section{Гарантии с доказательствами}
\label{sec:guarantees}

Межзапусковые утверждения используют \emph{инвариант повторного запуска}:
исправленный поток единиц, множества источников, граф, область действия и
глубина, базовые метки целей и разрешённый допуск каналов источников и целей
для каждой ожидающей единицы не меняются, кроме базовых меток успешно
выпущенных пакетов. Два компонента имеют имена. \textbf{(H1)} канал базовой
метки, читаемый \code{Propagate-Channel: inherit}, не меняется.
\textbf{(H2)} вклад, изначально подавленный отменой, удержанием или
исправлением, позднее не выпускается в той же $h$. Они необходимы в своих
случаях, но вместе не исчерпывают инвариант. Для детерминизма одного запуска
они не нужны.

\begin{lemma}[антитонность окна]
\label{lem:w}
Для $\sigma' = \mathit{run}(\sigma,F)$: $W^\circ_{\sigma'}(p)=\emptyset$,
если $p$ опубликован, и $W^\circ_{\sigma'}(p)=W^\circ_\sigma(p)$ иначе.
Окно серии $W$ также пустеет после стабильного выпуска, но может не
измениться после предрелизного.
\end{lemma}
\begin{proof}
Запуск добавляет только метки в $h$. Без новой метки обе базовые фиксации
неизменны. После публикации $\mathit{basec}_{\sigma'}(p)=h$, поэтому
формула~\ref{eq:fresh} вычитает $\mathit{reach}(h)$ и даёт пустое множество.
Только стабильный выпуск также делает $\mathit{scom}_{\sigma'}(p)=h$.
\end{proof}

\begin{lemma}[стабильность источников]
\label{lem:s}
При инварианте повторного запуска $\mathit{src}_{\sigma_j}(u)$ постоянно
для каждой ожидающей единицы на всей допустимой последовательности запусков.
\end{lemma}
\begin{proof}
Исправленный поток единиц и его предикаты отмены, удержания, исправления и
источников закреплены инвариантом. H2 исключает поздний выпуск изначально
подавленного вклада. Произвольное включение множеств меток вне допустимой
последовательности такой монотонности не даёт: метка с большей версией может
указывать на более раннюю фиксацию.
\end{proof}

\begin{lemma}[стабильность повторного запуска]
\label{lem:c}
При инварианте повторного запуска разрешённый канал невыпущенного пакета,
множество каналов источников каждой ожидающей единицы и
$\mathit{rsv}(u,d)$ постоянны.
\end{lemma}
\begin{proof}
Именно эти компоненты канала и множества источников закрепляет инвариант.
Определение~\ref{def:chan} выбирает то же прямое или распространённое
предложение, а формула~\ref{eq:resolvable} читает те же множества.
\end{proof}

\begin{remarknum}
\label{rem:drift}
В историческом утверждении~1.0 часть (H1) о базовой метке отсутствовала.
При одном лишь условии «все выпуски
стабильны» \code{inherit} читает
$\mathit{ch}(\mathit{base}_\sigma(q))$, которое сдвигает собственный
выход источника в стабильный канал: при \code{core} с базовой меткой \code{1.5.0-beta.0} и
фиксациях \code{feat(core)\^{}++1:~x} и
\code{release(core)\%stable:~graduate} запуск 1 планирует потребителя на
\code{2.0.1-beta.0} (унаследованный \code{beta}), пока \code{core}
выходит в стабильный канал; если потребитель отказывает, запуск 2 планирует его на \code{2.0.1}:
тот же пакет, оба плана, разные версии, каждый завершённый выпуск
стабилен. Этот исторический контрпример объясняет, почему CCME~2.0 включает
H1 в полный инвариант и не утверждает безусловную стабильность версий.
\end{remarknum}

\begin{theorem}[восемь гарантий]
\label{thm:g}
Обозначим $\mathit{adm}_j$ вместо $\mathit{adm}_{\sigma_j}$ и $n =
|\Pi(\sigma_1)|$. При инварианте повторного запуска для каждой единицы $u$, потребителя
$D$ и запуска $j \geq 1$:
\begin{description}
  \item[\G{1} завершаемость] распространение останавливается после $O(|\mathcal{U}|
    \cdot (|\mathcal{P}|{+}|E|))$ работы, каковы бы ни были объявленные глубины.
  \item[\G{2} отсутствие осиротевших пакетов] $\mathit{adm}_1(u,D) \Rightarrow
    \mathit{adm}_j(u,D)$, пока $D$ не выпущен в фиксации, содержащей
    $c_u$.
  \item[\G{3} стабильность версий] если $D \in \Pi(\sigma_1)$ и $D$ не
    выпущен до запуска $k$, то $D \in \Pi(\sigma_k)$ и $\nu_k(D) =
    \nu_1(D)$.
  \item[\G{4} отсутствие повторной доставки в журнале] как только $D$ выпущен после $c_u$,
    $\mathit{adm}_j(u,D)$ больше никогда не выполняется; вместе с \G{2} доставка
    в журнале происходит ровно один раз для каждого вклада, цель которого
    в конце концов успешно выпущена. Это не гарантирует однократность
    внешнего побочного эффекта.
  \item[\G{5} нерасширение] $\{D : \mathit{adm}_j(u,D)\} \subseteq
    \{D : \mathit{adm}_1(u,D)\}$.
  \item[\G{6} сходимость] $\Pi(\sigma_{n+1}) = \emptyset$ при условии, что
    каждая непустая повторная попытка допускает и публикует хотя бы один
    изначально публикуемый неудержанный пакет.
  \item[\G{7} погашение канала] распространённый вклад канала допускается,
    только пока его фиксация остаётся в $W^\circ(D)$; равенство базовому
    каналу и сопоставление перехода независимо отбрасывают пустые предложения.
  \item[\G{8} односторонняя независимость осей] подавление инкремента
    никогда не подавляет смену канала. Обратное ложно: разрешённые каналы
    входят в допуск инкремента в формуле~\ref{eq:resolvable}.
\end{description}
\end{theorem}

\begin{proof}
\G{1}: обход, вычисляющий $R$, помечает каждый увиденный пакет не более
одного раза, поэтому выполняет $O(|\mathcal{P}|{+}|E|)$ работы на единицу независимо от
глубины и циклов; ацикличность нужна для порядка публикации, а не здесь.
\G{2}: сравним формулу~\ref{eq:admission} условие за условием. Множество
источников, глубина, граф, область действия, отмена и допуск канала закреплены
инвариантом. Свежее окно потребителя постоянно до выпуска $D$
(лемма~\ref{lem:w}); только этот выпуск может сделать условие допуска ложным.
\G{3}: $D$ не получил метки, поэтому $\mathit{sbase}$, окно и
$\mathit{dir}$ не меняются; объединение распространённых инкрементов не меняется по \G{2}
($\geq$) и \G{5} ($\leq$), доказательства которых ниже используют только леммы; поэтому
$\mathit{eff}$ и канал не меняются (лемма~\ref{lem:c}), поэтому сохраняется
членство как по инкременту, так и только по каналу. Точное \code{Release-As}
остаётся тем же, поскольку его единица закреплена инвариантом. Во всех случаях
$\nu$ есть та же функция от неизменных входов.
\G{4}: по лемме~\ref{lem:w} $W^\circ(D)$ пустеет после стабильной или
предрелизной метки, исключая свежий прямой и распространённый допуск; также
$\mathit{pin}_\sigma(D)$ становится неопределённым. Это
погашение в журнале меток; разрыв между внешней публикацией и меткой требует
сверки тождества содержимого, а не гарантии однократности внешней операции.
\G{5}: каждое условие формулы~\ref{eq:admission} постоянно, кроме
окна, которое антитонно (лемма~\ref{lem:w}); конъюнкция
постоянных и антитонных предикатов антитонна.
\G{6}: $\Pi(\sigma_{j+1}) \subseteq \Pi(\sigma_j)$, поскольку член
более позднего плана не был выпущен (иначе \G{4}), и каждый дизъюнкт
$\Pi$ тогда не меняется или уменьшается; по предпосылке каждая непустая
попытка удаляет хотя бы один изначально публикуемый пакет (\G{4}), поэтому
$|\Pi|$ строго убывает до $0$ не более чем за
$n$ запусков. Запуск, который ничего не публикует, отказал не \emph{частично}
и выводится как фатальная ошибка, а не исполняется повторно.
\G{7}: проверка допуска оси каналов требует $c_u\in W^\circ(D)$, поэтому
лемма~\ref{lem:w} погашает вклад. Пока он свеж, равенство базовому каналу и
проверка \code{<from>} независимо отбрасывают пустые предложения.
\G{8}: проход каналов выполняется первым и не читает результат оси
инкрементов; следующий проход читает разрешённые каналы через
формулу~\ref{eq:resolvable}.
\end{proof}

\begin{proposition}[всеобщая расстановка меток необходима]
\label{prop:tagall}
Если бы выпущенный пакет без артефакта для загрузки и с ненулевым итоговым
инкрементом версионировался, но не помечался меткой, его окно никогда не
продвигалось бы (лемма~\ref{lem:w}), он вечно входил бы в каждый план заново,
и \G{6} нарушалась бы. Поэтому
определение~\ref{def:run} помечает меткой всё, что выпущено, а
проверка сходимости (повторное планирование после полностью успешного запуска; непустой
результат есть ошибка) допускает ровно одно исключение: удержанные пакеты.
\end{proposition}

Три следствия несут основную эксплуатационную нагрузку. \emph{Возобновление}: после
частичного отказа $\Pi(\sigma_2) = F \cup B$ на версиях запуска 1
(\G{2}--\G{5}), поэтому следующий запуск доделывает остаток той же работы без
флагов и без очистки. \emph{Идемпотентность}: два запуска без промежуточных
фиксаций и изменений меток или конфигурации дают один и тот же план
(теорема~\ref{thm:det});
запуск после полного успеха даёт пустой план (\G{4}).
\emph{Корректность рецензирования}: пока действует инвариант повторного
запуска, по \G{5} план, который рецензент одобряет перед
первой попыткой, есть верхняя граница всего, что может опубликовать любая повторная попытка,
поэтому злоумышленник, способный вызвать отказ публикации, не может расширить
то, что выпускает фиксация, и ни одно возобновление не требует повторного одобрения.

\begin{table}[t]
\caption{Раздел~\ref{sec:motivation} при расписании, которое дважды теряет
\code{cli}. \G{2}: \code{cli} входит в каждый план, пока не
опубликуется. \G{3}: во всех планах на \code{2.0.1}. \G{4}: три опубликованных
пакета покидают план. \G{5}: ничего нового не входит. \G{6}: план становится
пустым. Запуск 2 ничего не публикует, что \G{6} относит к случаю фатальной
ошибки, а не частичного отказа.}
\label{tab:runs}
\small
\begin{tabularx}{\columnwidth}{@{}lYYY@{}}
\toprule
Запуск & План & Опубликовано & Отказ \\
\midrule
1 & \code{core}, \code{cli}, \code{ui}, \code{api}
  & \code{core} 1.5.0, \code{ui} 0.9.2, \code{api} 1.2.1 & \code{cli} \\
2 & \code{cli} 2.0.1 & --- & \code{cli} \\
3 & \code{cli} 2.0.1 & \code{cli} 2.0.1 & --- \\
4 & $\emptyset$ & --- & --- \\
\bottomrule
\end{tabularx}
\end{table}

\section{Конкурирующие средства как конкретизации модели}
\label{sec:compare}

В исполненных зафиксированных конфигурациях Lerna и nx выводят инкременты из
коммитов и ставят метки до публикации, поэтому отказавшая загрузка отсутствует
в их журнале меток; документация release-please задаёт тот же порядок.
Changesets ставит метки после публикации, поглощает файлы намерений и при
восстановлении спрашивает npm. CCME помечает каждый пакет после его публикации
и сохраняет долг в git с ограничением сверки внешней публикации и метки из
определения~\ref{def:run}. Наблюдавшееся распространение было безусловным у
Lerna/nx, ограниченным диапазонами у Changesets и задаваемым каждым коммитом у
CCME. Полные протоколы команд сохранены в артефакте.

Фикстуры отказа закрепляют Lerna~10.0.1, Nx~23.1.2 и Changesets~3.0.1.
Lerna/nx поставили метку отклонённой версии до неудачной загрузки; Changesets
оставил её без метки и восстановился запросом к npm. При параллельной передаче
сравниваемые запуски оставили метки вне предков ветви, а dispat слил ветви до
последней передачи и запланировал коммит коллеги следующим.

\begin{table}[t]
\caption{Зафиксированные конфигурации средств выпуска.}
\label{tab:compare}
\scriptsize
\resizebox{\columnwidth}{!}{%
\begin{tabular}{@{}llll@{}}
\toprule
 & lerna / nx & changesets & CCME / dispat \\
\midrule
источник & коммиты & файлы намерений & коммиты \\
время метки & до публикации & после публикации & после каждого шага \\
распространение & безусловное & по диапазонам & по коммиту \\
журнал восстановления & метки / CI & npm & метки git \\
внешний разрыв & вручную/запрос & запрос npm & сверка тождества \\
\bottomrule
\end{tabular}}
\end{table}
\section{Реализация}
\label{sec:impl}

dispat --- статически скомпонованный исполняемый файл Go, чьё состояние
планирования выпуска состоит только из конфигурации, истории и меток. Пакеты
являются каталогами, а этапы жизненного цикла --- командами оболочки;
сканеры и сохраняющие байты средства записи поддерживают 35 форматов
манифестов двадцати экосистем. Планировщик реализует данные выше определения;
топологическое исполнение продолжает независимые поддеревья после отказа и
пропускает потребителей отказавших пакетов. Удалённая зарезервированная метка
сериализует запуски.

Успешные внешние эффекты не откатываются: последующий отказ метки, записи или
передачи оставляет работу по сверке согласно \S19.4 CCME~2.0. Вызываемые
средства, включая Terraform, отвечают за своё операционное состояние и
проверку тождества. Реализация проверяется векторами, чёрноящичными запусками
CLI, фаззингом, реальными манифестами и описанными ниже кампаниями.

\paragraph{Структуры репозиториев и гранулярность саги.}
Для репозитория с одним пакетом $\mathcal{P}=\{p\}$ и $E=\varnothing$:
распространение пусто, но планирование версии, запись «публикация, затем метка»
и восстановление сохраняются. Это одношаговый случай протокола саги. Dispat
требует поместить пакет в подкаталог; конфигурация и журнал выпуска остаются в
корне. Несколько назначений внутри одного сценария публикации разделяют одну
метку завершения, поэтому частичный внешний успех требует сверки; отдельные
пакеты-шаги дают отдельные записи завершения. Один исходный проект может
представить такими шагами CLI, образ контейнера и документацию.

Для полирепозитория управляющий репозиторий закрепляет связанные репозитории
как подмодули git внутри каталогов-обёрток пакетов. Его коммиты указателей
выражают намерение выпуска, а его история и метки дают единственные $h$ и
$\sigma$ модели. Связанные репозитории сохраняют собственные истории. Так
получается единый граф выпуска без слияния исходных репозиториев, а не
распределённая транзакция над независимыми историями git. Конструкция
показывает применимость; оценка не измеряет эксплуатационный масштаб
полирепозитория.

\paragraph{Повторно используемая автоматизация.}
Разработчик в Windows и раннер CI на Linux могут использовать общий
\code{global.yaml} с графом, сценариями и жизненным циклом. Каждый платформенный
файл ссылается на него через \code{\$ref: ./global.yaml}; \code{linux.yaml}
задаёт \code{shell: ["/bin/sh", "-c"]}, а \code{windows.yaml} задаёт оболочку
\code{["powershell.exe", "-NoProfile", "-Command"]}.
Повторяются только платформенные переопределения; ссылки относительны к файлу,
ключи заменяются целиком, а не по вложенным полям. Локально план просматривают
через \code{dispat --config windows.yaml status}, на раннере --- через
\code{dispat --config linux.yaml status}; после барьеров выпуска CI запускает
последнюю конфигурацию без \code{status}.
\code{dispat if OS=Windows\_NT} выбирает эти команды через
\code{--then} и \code{--else}, но версия~1.8.2 исполняет ветви через
\code{/bin/sh} до чтения конфигурации дочерним вызовом. В нативной Windows
можно выбрать файл напрямую; условная обёртка всё ещё требует оболочки POSIX.
Именованные сценарии также могут использовать \code{\$ref}; команда
\code{dispat for Debug Release --do "dispat exec build-one"}
передаёт каждое значение через \code{DISPAT\_ITEM}.

\paragraph{Барьеры выпуска для агента.} Авторизация отдельна. Проверка
\code{status} с флагами \code{--require-release} и \code{--strict}
без блокировки может
предшествовать набору тестов всего рабочего пространства с проверенной
очисткой временных ссылок. Затем проверки пакетов подтверждают согласованные
версии и lock-файлы, а дымовые тесты --- точные артефакты перед публикацией.
Обёртки сохраняют отказы.
PR~2813 HOL Guard добавляет включаемое рецензирование начала выпуска
Dispat~\cite{guard}; тестовые барьеры предоставляют Dispat и сценарии проекта.
На 6 сентября 2026 года это предложение интеграции оставалось открытым, а его
полный набор тестов --- непроверенным. Это пример применения реализации, а не
развёртывание или дополнительная гарантия протокола.
\section{Оценка}
\label{sec:eval}

Если явно не указана свежая кампания CCME~2.0/dispat~1.8.2, ниже сохранены
исторические измерения названных образов~1.7.x. Артефакт фиксирует образы,
начальные числа, сырые потоки, контрольные суммы и версии стендов.

\subsection{Дифференциальная проверка уточнения}
\label{sec:eval:diff}

Открытым в разделе~\ref{sec:model:inst} остаётся соответствие бинарного файла
модели. Оракул примерно из 300 строк Python реализует только
раздел~\ref{sec:model}; генератор создаёт 5--20 пакетов, слоистые DAG и 4--12
единиц, реализует их настоящим репозиторием git и сравнивает план с выпущенным
образом dispat~1.7.2. Золотые случаи из статьи служат предварительным барьером.

В области старых (H1),(H2) ядро исторического оракула~1.0 совпало точно на
1{,}500/1{,}500 экземплярах. Полный профиль намеренно включает неполные в
модели~1.0 предрелизные базы, записи только канала и каре у типов без
инкремента. На 1{,}000 экземплярах он дал 1{,}097 различающихся строк пакетов,
все в ожидаемых классах: 409 входов в канал, 240 погашений серии, 58 вариантов
утечки подавления и каскада и 390 случаев живости оси инкремента. Старое
правило допуска опускало $\mathit{live}$: не требовало ненулевых собственного и
распространённого инкрементов и положительной глубины. Ещё 500 начальных чисел
дали те же классы без неклассифицированных строк.

\subsection{Свежая кампания CCME 2.0 / dispat 1.8.2}
\label{sec:eval:fresh}

Выпущенная CCME~2.0.0 повторно проверена официальным ARM64-образом
Dispat~1.8.2. Кампания связывает коммит тега \code{a7011cbcb9e8}, дайджест
образа \code{sha256:4c0ad4b8...e7bfd}, проверенный файл выпуска и 753 артефакта.

Стенд проверяет фикстуры Git и коды CLI, использует tmpfs для E1 и отделяет
семантические отказы от ошибок исполнения. Сохранённые записи 1.8.1 получены
прежним стендом; улучшение реализации из сравнения не выводится. Все 1{,}500 экземпляров core совпали,
непригодных нет. Из 1{,}500 историй полного профиля, включающего
регрессивные предварительные версии, 997
дали принятые планы (546 чистых сравнений), а 503 --- полные диагностические
отказы E195; четыре отказа не дали различий, итого диагностически чисты 550.
Для всех 503 досье независимо проверены предложенные регрессивные версии.
Остальные 1{,}914 строк имеют классы PG1 862, PG2 375, PG2x 14, PG3? 663 и
UNKNOWN 0. Отклонённые диагностические планы не являются исполнимыми планами
или результатами проверки соответствия. CLI прошёл выбранные векторы 133--138 версии~2.0: свежее
погашение, подавление, допуск и порядок каналов, дрейф наследуемого канала и
исключение источника. Вектор~135 наблюдает только нисходящий W208, не
восходящий API; шесть случаев не доказывают уточнение.

Два повторения осиротевшего пакета восстановили ту же версию и пустой третий
план; две кампании сошлись в 24/24 эпизодах за не более чем четыре запуска с
нулём структурных ошибок G3/G5. Воспроизведение пяти выпусков получило 36/42
версий google-cloud-node и 7/64 conventional-changelog, не утверждая, что те
выпуски использовали CCME. Три свежие ячейки dispat прошли 8, 12 и 14 проверок; конкуренты
повторно не запускались.

Все восемь ячеек времени --- допустимые наблюдения на нагруженном узле при
ограничении 2 CPU/4 ГиБ, а не межинструментальное сравнение. Медианы лежат от
0.017\,с для 100 слоистых пакетов до 19.354\,с для цепочки из 10{,}000 пакетов;
наблюдаемый максимум RSS --- 15--118\,МиБ. Для U1000 получено
0.290\,с/26\,МиБ, для 5{,}000 дополняющих коммитов --- 0.664\,с/30\,МиБ.
Записанные фикстуры прошли \code{git fsck}; генератор \code{commit-tree}
совпал с обычной малой фикстурой по идентификаторам коммитов, ссылкам,
сообщениям и журналу.

\subsection{Способность формализации находить дефекты}
\label{sec:eval:formal}

Формализация исторической CCME~1.0.0 выявила шесть мест, где утверждение не
следовало из правил. CCME~2.0.0 включает исправления; контрпримеры остаются
регрессионными обязательствами, а не дефектами выпущенной версии~2.0.
\begin{enumerate}
\item Окно серии $W$ сохраняет накопление для формулы~\ref{eq:train}, а допуск
и погашение используют свежее $W^\circ$ из формулы~\ref{eq:fresh}; иначе
повторяется \code{beta.1}, \code{beta.2},~\ldots
\item Истекающее подавление может расширить \G{5}; версия~2.0 включает H2 в
полный инвариант и требует рецензирования при его изменении.
\item Исторической \G{3} не хватало стабильности наследуемого канала:
\code{inherit} читает базу источника, которую может изменить выпуск в stable.
Версия~2.0 включает это условие H1 в инвариант.
\item Старый восходящий аудит опускал формулу~\ref{eq:resolvable}; версия~2.0
применяет полный нисходящий предикат в обратном направлении, включая допуск
канала и исключение самой цели из источников.
\item Симметричная «независимость осей» была ложна: отвергнутый канал может
затем подавить инкремент. \G{8} версии~2.0 односторонняя, каналы разрешаются
первыми.
\item Объединение корзин теряло исключение источника отдельной единицы.
Версия~2.0 повторяет эту проверку для каждой единицы.
\end{enumerate}

Наконец, найдено текущее текстовое несогласование: явное правило \S11.5
правильно задаёт для выпуска из серии совокупность $E$, тогда как резюме
\S13.9 говорит о свежем $\mathit{eff}$ и должно ссылаться на \S11.5. E185
требует цель строго выше ядра базы, хотя нормативный пример выпускает
\code{1.3.0-rc.1} в равное ядро \code{1.3.0}. Это несогласованности ссылки и
защиты, не утверждение о дефекте исполнения.
\section{Угрозы обоснованности}
\label{sec:threats}

\emph{Внутренняя.} Векторы, тесты и модель созданы автором системы и могут
разделять ошибку. Гарантии относятся к модели; дифференциальная проверка
раздела~\ref{sec:eval:diff} связывает её с планировщиком, но не с протоколом
запуска. Эксперимент параллельной передачи --- сильнейшее свидетельство для
протокола, однако это одно детерминированное чередование на средство.
\emph{Внешняя.} Общность по экосистемам опирается на конструкцию, проверенные
средства записи и выбранные нами фикстуры, а не на полевое исследование.
\emph{Конструктная.} Контракт развивался вместе с реализацией;
раздел~\ref{sec:eval:formal} показывает обратную связь. Рецензируемость не
проверена пользователями: утверждается лишь запись области воздействия в
сообщении. Release-please характеризуется по документации; Lerna, nx и
Changesets исполнены на двух отказах в закреплённой шестипакетной фикстуре из
раздела~\ref{sec:compare}. Их проверки описывают оставшееся состояние, а не
общую оценку. Контейнерные эксперименты используют одну машину и одного
поставщика CI; восстановленные конфигурации репозиториев могли неверно понять
их соглашения.

\section{Смежные работы}
\label{sec:related}

Инженерия выпусков~\cite{adams} и монорепозиторий Google~\cite{potvin}
изучены ранее. Semantic-release
\cite{semanticrelease} развил версионирование по коммитам; Lerna~\cite{lerna},
Changesets~\cite{changesets}, Nx~\cite{nx} и release-please
\cite{releaseplease} расширили автоматизацию на монорепозитории. Эмпирические
работы показывают расхождение объявленной семантики версий и совместимости
\cite{raemaekers,bogart}, а Decan и др. изучают эволюцию сетей зависимостей
\cite{decan}. Формальные модели операций VCS~\cite{swierstraloh}, теорий
патчей~\cite{hpt} и слияния~\cite{mimram} описывают содержимое; здесь остаются
DAG, метки и сообщения, используемые анализом git~\cite{birdgit} и моделями
слияния~\cite{mrdt}. Пакеты и графы восходят к EDOS~\cite{edos}; в пределах
рассмотренной литературы их ранее не объединяли в данный протокол выпуска.

Протокол адаптирует саги~\cite{sagas} к шагам без приемлемой компенсации и
выбирает продвижение вперёд в границах транзакций Грея~\cite{gray}.
Независимая фиксация и повторение Хелланда~\cite{helland07,helland12}
соответствуют \G{2}--\G{6}. Порядок публикации
основан на алгоритме Кана~\cite{kahn}; таксономия систем сборки
\cite{mokhov} отделяет отсутствие неизменённых пакетов в плане от кэша
пересборки. Фаззинг учитывает предостережения Клис и др.~\cite{klees}. Anvil
\cite{anvil} доказывает живость Kubernetes-контроллеров; здесь сходимость дают
журнал, повтор и сверка неопределённых внешних исходов.

\section{Заключение}
\label{sec:conclusion}

Для репозиториев с одним пакетом, монорепозиториев и полирепозиториев с
управляющим репозиторием CCME рассматривает публикацию выпуска как сагу.
Он пишет намерение в сообщениях, завершённые шаги --- в метках и допускает
распространённую работу по свежему окну потребителя. При инварианте повтора
восемь гарантий делают детерминированный повтор корректным восстановлением.
Шесть дефектов исторического текста~1.0 исправлены в CCME~2.0; свежие векторы,
восстановление и закреплённые отказы подтверждают заявленные границы.
Механизация в Lean или TLA$^{+}$~\cite{lamport} и случайные чередования полных
запусков остаются будущей работой.

\section*{Доступность данных}

Выпущенные исходный код и спецификации доступны по адресу
\url{https://github.com/yohimik/dispat}. Локальный сопроводительный набор
хранит детали доказательств, векторы, фикстуры, стенды, контрольные суммы и
сырые записи; его архивное размещение остаётся требованием готовности к
публикации.

\vfuzz=2pt
\begin{thebibliography}{21}

\bibitem{sagas}
H.~Garcia-Molina and K.~Salem.
Sagas.
In \emph{Proceedings of the ACM SIGMOD International Conference on
Management of Data}, pages 249--259, 1987.

\bibitem{gray}
J.~Gray.
The transaction concept: Virtues and limitations.
In \emph{Proceedings of the 7th International Conference on Very Large
Data Bases (VLDB)}, pages 144--154, 1981.

\bibitem{helland07}
P.~Helland.
Life beyond distributed transactions: an apostate's opinion.
In \emph{Proceedings of the 3rd Biennial Conference on Innovative Data
Systems Research (CIDR)}, 2007.

\bibitem{helland12}
P.~Helland.
Idempotence is not a medical condition.
\emph{Communications of the ACM}, 55(5):56--65, 2012.

\bibitem{kahn}
A.~B. Kahn.
Topological sorting of large networks.
\emph{Communications of the ACM}, 5(11):558--562, 1962.

\bibitem{semver}
T.~Preston-Werner.
Semantic Versioning 2.0.0.
\url{https://semver.org/spec/v2.0.0.html}, accessed 2026.

\bibitem{cc}
Conventional Commits 1.0.0.
\url{https://www.conventionalcommits.org/en/v1.0.0/}, accessed 2026.

\bibitem{ccme2}
CCME Specification 2.0.0 and historical 1.0.0 snapshot.
\url{https://github.com/yohimik/dispat/blob/d3ce8ae637232908ced82efb5cda550f60ed2fdf/specs/ccme-spec/SPEC.md};
\url{https://github.com/yohimik/dispat/blob/713f1a0b/pkg/ccme/SPEC.md},
accessed 2026.

\bibitem{ccme301}
CCME Specification 3.0.1, corrected algorithm guidance.
\url{https://github.com/yohimik/dispat/blob/8238e2465670b39609f8dcf5a8a5d34ea71d3db1/specs/ccme-spec/SPEC.md},
accessed 2026.

\bibitem{birdgit}
C.~Bird, P.~C. Rigby, E.~T. Barr, D.~J. Hamilton, D.~M. Germ\'an, and
P.~Devanbu.
The promises and perils of mining Git.
In \emph{Proceedings of the 6th IEEE International Working Conference
on Mining Software Repositories (MSR)}, 2009.

\bibitem{mrdt}
G.~Kaki, S.~Priya, K.~C. Sivaramakrishnan, and S.~Jagannathan.
Mergeable replicated data types.
\emph{Proceedings of the ACM on Programming Languages},
3(OOPSLA), Article 154, 2019.

\bibitem{swierstraloh}
W.~Swierstra and A.~L\"oh.
The semantics of version control.
In \emph{Proceedings of the ACM International Symposium on New Ideas,
New Paradigms, and Reflections on Programming \& Software (Onward!)},
pages 43--54, 2014.

\bibitem{hpt}
C.~Angiuli, E.~Morehouse, D.~R. Licata, and R.~Harper.
Homotopical patch theory.
In \emph{Proceedings of the 19th ACM SIGPLAN International Conference
on Functional Programming (ICFP)}, 2014.

\bibitem{mimram}
S.~Mimram and C.~Di Giusto.
A categorical theory of patches.
\emph{Electronic Notes in Theoretical Computer Science}, 298, 2013.

\bibitem{edos}
F.~Mancinelli, J.~Boender, R.~Di Cosmo, J.~Vouillon, B.~Durak,
X.~Leroy, and R.~Treinen.
Managing the complexity of large free and open source package-based
software distributions.
In \emph{Proceedings of the 21st IEEE/ACM International Conference on
Automated Software Engineering (ASE)}, 2006.

\bibitem{lerna}
Lerna: a fast, modern build system for managing and publishing multiple
JavaScript packages from the same repository.
\url{https://lerna.js.org/}, accessed 2026.

\bibitem{changesets}
Changesets: a way to manage your versioning and changelogs.
\url{https://github.com/changesets/changesets}, accessed 2026.

\bibitem{semanticrelease}
semantic-release: fully automated version management and package
publishing.
\url{https://github.com/semantic-release/semantic-release},
accessed 2026.

\bibitem{nx}
Nx: smart monorepos, fast CI.
\url{https://nx.dev/}, accessed 2026.

\bibitem{releaseplease}
release-please: automated releases based on conventional commits.
\url{https://github.com/googleapis/release-please}, accessed 2026.

\bibitem{adams}
B.~Adams and S.~McIntosh.
Modern release engineering in a nutshell: Why researchers should care.
In \emph{Proceedings of the 23rd IEEE International Conference on
Software Analysis, Evolution, and Reengineering (SANER)}, 2016.

\bibitem{potvin}
R.~Potvin and J.~Levenberg.
Why Google stores billions of lines of code in a single repository.
\emph{Communications of the ACM}, 59(7):78--87, 2016.

\bibitem{raemaekers}
S.~Raemaekers, A.~van Deursen, and J.~Visser.
Semantic versioning versus breaking changes: A study of the Maven
repository.
In \emph{Proceedings of the 14th IEEE International Working Conference
on Source Code Analysis and Manipulation (SCAM)}, 2014.

\bibitem{bogart}
C.~Bogart, C.~K\"astner, J.~Herbsleb, and F.~Thung.
How to break an API: Cost negotiation and community values in three
software ecosystems.
In \emph{Proceedings of the 24th ACM SIGSOFT International Symposium on
Foundations of Software Engineering (FSE)}, 2016.

\bibitem{decan}
A.~Decan, T.~Mens, and P.~Grosjean.
An empirical comparison of dependency network evolution in seven
software packaging ecosystems.
\emph{Empirical Software Engineering}, 24(1):381--416, 2019.

\bibitem{mokhov}
A.~Mokhov, N.~Mitchell, and S.~Peyton Jones.
Build systems \`a la carte.
\emph{Proceedings of the ACM on Programming Languages},
2(ICFP):79:1--79:29, 2018.

\bibitem{klees}
G.~Klees, A.~Ruef, B.~Cooper, S.~Wei, and M.~Hicks.
Evaluating fuzz testing.
In \emph{Proceedings of the ACM SIGSAC Conference on Computer and
Communications Security (CCS)}, pages 2123--2138, 2018.

\bibitem{anvil}
X.~Sun, W.~Ma, J.~T. Gu, Z.~Ma, T.~Chajed, J.~Howell, A.~Lattuada,
O.~Padon, L.~Suresh, A.~Szekeres, and T.~Xu.
Anvil: Verifying liveness of cluster management controllers.
In \emph{Proceedings of the 18th USENIX Symposium on Operating Systems
Design and Implementation (OSDI)}, 2024.

\bibitem{lamport}
L.~Lamport.
\emph{Specifying Systems: The TLA\textsuperscript{+} Language and
Tools for Hardware and Software Engineers}.
Addison-Wesley, 2002.

\bibitem{guard}
Semyon F. (yohimik). Opt-in Dispat release-start review.
\emph{hashgraph-online/hol-guard}, pull request 2813, 2026.
\url{https://github.com/hashgraph-online/hol-guard/pull/2813}.
Head \texttt{132bb5d105e1}; accessed 6 September 2026.

\end{thebibliography}

\end{document}
